\begin{theorem}[Church, 1936]
    Każda funkcja rekurencyjna ma \(\lambda\)-reprezentację.
\end{theorem}

\begin{theorem}
    Każda funkcja \(\lambda\)-reprezentowalna jest rekurencyjna.
\end{theorem}
\section{Arytmetyzacja rachunku lambda}
Enumerujemy symbole:
\begin{center}
\begin{tabular}{ r | c c c c c c c }
 znak \(x_i\) & \(\cdot\) & ( & ) & \(\lambda\) & \(x_1\) & \(x_2\) & \dots \\ 
kod \(\alpha_i\) & 2 & 3 & 5 & 7 & 11 & 13 & \dots
\end{tabular}
\end{center}

Kodujemy skończone typy: \\
Ciąg \( x_1 \ldots x_n \) kodujemy jako \( p_1^{\alpha_1} \ldots p_n^{\alpha_n} \). \\

Definiujemy \(q, r \in \prec\) dla termów \(S, T\) w postaci normalnej kodowanych przez \(n, k\): \\
\(q(n,k)\) -- numer redukcji \(T \to_{\alpha\beta} S\) \\
\(r(n,k)\) -- numer równości \(T =_{\alpha\beta} S\) \\
Definicje są poprawne i \(q,r\) rzeczywiście są rozstrzygalne, ponieważ możemy ponumerować wszystkie postaci normalne. \\

Konstruujemy rodzinę predykatów \( \tau^n \subseteq \natural^{n+2} \):
\[ \tau^n(e, x_1, \ldots, x_n, y), \]
oznacza, że \(y\) jest numerem redukcji termu o numerze \(e\) do termu \(x_1, \ldots, x_n\). \\

Konstruujemy funkcję \(U: \natural \to \natural\):
\begin{align*}
    & U(y) = n, \text{ gdy } y \text{ to numer redukcji do } n \\
    & f(x_1, \ldots, x_n) = U(n, y, \tau^n(e, x_1, \ldots, x_n), y)
\end{align*}